% ==============================================================================
% BÁO CÁO KHOA HỌC: HỆ THỐNG PHÁT HIỆN Ổ GÀ VÀ VẬT CẢN ĐƯỜNG BỘ VÀ CẢNH BÁO SỚM
% THỜI GIAN THỰC DỰA TRÊN MÔ HÌNH YOLO11n
% Tác giả: Hoàng Quang Minh, Nguyễn Đức Tâm & Lê Văn Tịnh
% 
% File tổng hợp bao gồm 3 phần trọng tâm:
%   MỤC 3: PHƯƠNG PHÁP NGHIÊN CỨU (METHODOLOGY)
%   MỤC 4: TRIỂN KHAI PHẦN MỀM (SOFTWARE DEPLOYMENT)
%   MỤC 5: KẾT QUẢ VÀ THẢO LUẬN (RESULTS AND DISCUSSION)
% Format chuẩn hóa theo phong cách Tech Science Press (CMC)
% ==============================================================================

% ==============================================================================
% MỤC 3: PHƯƠNG PHÁP NGHIÊN CỨU (METHODOLOGY)
% ==============================================================================
\section{PHƯƠNG PHÁP NGHIÊN CỨU (METHODOLOGY)}
\label{sec:methodology}

Để giải quyết đồng thời bài toán nhận diện chính xác các hư hỏng mặt đường dạng âm (ổ gà - \textit{potholes}), các chướng ngại vật kích thước nhỏ, đa dạng và bảo đảm khả năng phản hồi thời gian thực trên các thiết bị tính toán biên có tài nguyên hạn chế, nghiên cứu đề xuất một khung phương pháp luận toàn diện gồm ba hợp phần:
\begin{enumerate}
    \item \textbf{Xây dựng và chuẩn hóa bộ dữ liệu đa lớp chuyên dụng (Constructing the Dataset):} Quy trình 4 bước thu thập, tiền xử lý, tăng cường dữ liệu và biên soạn bộ dữ liệu đường bộ gồm 15 lớp đối tượng vật cản đặc thù tại Việt Nam.
    \item \textbf{Kiến trúc mô hình YOLO11n (Network Architecture of YOLO11n):} Trình bày cấu trúc mạng học sâu nhẹ tiên tiến YOLO11n với các module C3k2, SPPF, đầu phát hiện tách rời không neo (Decoupled Anchor-free Head) và phân tích lý thuyết lựa chọn bộ tối ưu hóa SGDM thay vì Adam.
    \item \textbf{Cơ sở lý thuyết và các chỉ số hiệu năng (Theoretical Foundations and Performance Metrics):} Hợp nhất toàn bộ cơ sở lý thuyết, công thức toán học của bài toán phát hiện đối tượng, ước lượng khoảng cách độ sâu, động học đánh giá rủi ro va chạm và các chỉ số đo kiểm thời gian thực của phần mềm.
\end{enumerate}

% ------------------------------------------------------------------------------
\subsection{Xây Dựng Bộ Dữ Liệu (Constructing the Dataset)}
\label{subsec:xay_dung_du_lieu}

Quy trình chuẩn bị dữ liệu được thực hiện nghiêm ngặt qua 4 giai đoạn chuẩn hóa:
\begin{enumerate}
    \item \textbf{Thu thập dữ liệu (Data Collection):} Hình ảnh được tổng hợp có chọn lọc từ các kho dữ liệu công khai (Roboflow Universe Projects, Pothole-600, RTK Road Obstacle) kết hợp với các đoạn video hành trình (dashcam footage) ghi nhận thực tế trên đường phố Việt Nam ở độ phân giải cao. Dữ liệu bao quát đa dạng điều kiện môi trường: nắng gắt chói lóa, trời râm mát, bóng râm cây cối đổ trên mặt đường, hoàng hôn và ban đêm có ánh đèn pha.
    \item \textbf{Tiền xử lý dữ liệu (Preprocessing):} Loại bỏ thông tin hướng xoay EXIF (\textit{EXIF Orientation Tags}) nhằm ngăn chặn hiện tượng xoay ngược khung hình. Các nhãn đối tượng được chuẩn hóa về định dạng nhãn chuẩn YOLO (\texttt{.txt}). Ảnh đầu vào được đưa về kích thước vuông chuẩn $640 \times 640$ pixels bằng phương pháp căn đệm bảo toàn tỷ lệ khung hình (\textit{Letterbox Resizing}).
    \item \textbf{Tăng cường dữ liệu (Data Augmentation):} Áp dụng chuỗi kỹ thuật tăng cường đa dạng:
    \begin{itemize}
        \item \textit{Biến đổi không gian:} Lật ảnh đối xứng gương theo trục ngang (\textit{Horizontal Flip}) với xác suất $p = 0,5$; xoay ngẫu nhiên trong góc hẹp $\theta \in [-15^\circ, +15^\circ]$; thu phóng và cắt xén ngẫu nhiên (\textit{Random Crop \& Zoom}) trong biên độ $0\% - 20\%$.
        \item \textit{Biến đổi quang học và màu sắc:} Nhiễu màu ngẫu nhiên trong không gian HSV (Hue $\pm 0,015$, Saturation $\pm 0,7$, Value/Brightness $\pm 0,4$) nhằm mô phỏng các điều kiện phơi sáng và bóng tối cục bộ trên mặt đường.
        \item \textit{Ghép ảnh Mosaic:} Kết hợp 4 mảnh hình ảnh từ các ngữ cảnh khác nhau thành một khung hình tổng hợp với hệ số $mosaic = 1,0$ trong 90 epochs đầu tiên và tắt hoàn toàn (\textit{Close Mosaic}) trong 10 epochs cuối để ổn định gradient.
    \end{itemize}
    \item \textbf{Đóng gói và tổng hợp bộ dữ liệu (Compilation):} Toàn bộ dữ liệu sau xử lý được đóng gói thành bộ dữ liệu giao thông đường bộ hoàn chỉnh gồm 21.272 ảnh và 70.935 nhãn vật thể được phân định ranh giới chuẩn xác.
\end{enumerate}

Bảng~\ref{tab:image_database_sources} mô tả chi tiết 5 nguồn thành phần tạo nên cơ sở dữ liệu ảnh giao thông đường bộ phục vụ huấn luyện và kiểm thử.

\begin{table}[htbp]
\centering
\caption{Chi tiết các nguồn thành phần trong cơ sở dữ liệu ảnh đường bộ và vật cản.}
\label{tab:image_database_sources}
\small
\begin{tabular}{|l|c|c|c|p{5.8cm}|}
\hline
\textbf{Nhóm dữ liệu nguồn} & \textbf{Số lượng ảnh} & \textbf{Số lượng nhãn} & \textbf{Kích cỡ ảnh (MP)} & \textbf{Mô tả ngữ cảnh thu nhận} \\ \hline
Dữ liệu ổ gà và hư hỏng mặt đường & 3.850 & 7.225 (TB: 1,88) & 0,45 -- 2,10 & Ảnh ổ gà, nứt sụt lún trên mặt đường asphalt và bê tông dưới nhiều góc nhìn camera. \\ \hline
Dữ liệu phương tiện giao thông chính & 8.420 & 23.209 (TB: 2,76) & 1,20 -- 4,15 & Luồng ô tô, xe tải, xe buýt, xe máy trong điều kiện mật độ giao thông đô thị cao. \\ \hline
Dữ liệu người đi bộ và đối tượng yếu thế & 4.680 & 30.888 (TB: 6,60) & 0,85 -- 2,50 & Người đi bộ, xe đạp, động vật xuất hiện ngẫu nhiên băng qua lòng đường hoặc mép lề. \\ \hline
Dữ liệu rào chắn và phân làn công trình & 2.870 & 4.951 (TB: 1,73) & 0,50 -- 2,20 & Chóp nón phản quang, cọc tiêu cảnh báo, khối cầu chắn rào tại các khu vực đang sửa đường. \\ \hline
Dữ liệu vật thể ven đường và biển báo & 1.452 & 4.662 (TB: 3,21) & 0,60 -- 2,00 & Thùng rác, trụ cứu hỏa, biển báo dừng ven vỉa hè dễ gây cản trở tầm nhìn tài xế. \\ \hline
\textbf{Tổng cộng} & \textbf{21.272} & \textbf{70.935 (TB: 3,33)} & \textbf{0,45 -- 4,15} & \textbf{Bao quát toàn diện các tình huống giao thông thực tế.} \\ \hline
\end{tabular}
\end{table}

Mỗi tệp gán nhãn tương ứng với một ảnh được mã hóa dưới dạng bảng văn bản chuẩn, trong đó mỗi hàng đại diện cho một đối tượng với 5 trường thông số hình học được mô tả trong Bảng~\ref{tab:details_of_annotations}.

\begin{table}[htbp]
\centering
\caption{Đặc tả cấu trúc các trường thông tin trong tệp gán nhãn (Annotation Format).}
\label{tab:details_of_annotations}
\small
\begin{tabular}{|l|c|p{7.8cm}|}
\hline
\textbf{Trường dữ liệu} & \textbf{Miền giá trị / Định dạng} & \textbf{Mô tả ý nghĩa kỹ thuật} \\ \hline
\texttt{<class\_index>} & Số nguyên $[0, 14]$ & Chỉ số định danh lớp của đối tượng trong không gian 15 danh mục cần nhận diện. \\ \hline
\texttt{<x\_coordinate>} & Số thực $[0, 1]$ & Tọa độ tâm $x$ của hộp giới hạn, chuẩn hóa theo chiều rộng ảnh (0: mép trái, 1: mép phải). \\ \hline
\texttt{<y\_coordinate>} & Số thực $[0, 1]$ & Tọa độ tâm $y$ của hộp giới hạn, chuẩn hóa theo chiều cao ảnh (0: mép đỉnh, 1: mép đáy). \\ \hline
\texttt{<width>} & Số thực $[0, 1]$ & Chiều rộng tương đối của hộp giới hạn so với chiều rộng toàn khung hình. \\ \hline
\texttt{<height>} & Số thực $[0, 1]$ & Chiều cao tương đối của hộp giới hạn so với chiều cao toàn khung hình. \\ \hline
\end{tabular}
\end{table}

Toàn bộ 21.272 ảnh trong cơ sở dữ liệu được phân chia ngẫu nhiên nhưng phân tầng (\textit{Stratified Split}) thành 3 tập dữ liệu độc lập theo tỷ lệ tương ứng xấp xỉ $71,9\% : 16,4\% : 11,7\%$ (tương đương cấu hình $80:15:5$ tiêu chuẩn đối với tập dữ liệu huấn luyện mở rộng), chi tiết thể hiện tại Bảng~\ref{tab:dataset_split_statistics}.

\begin{table}[htbp]
\centering
\caption{Phân chia bộ dữ liệu phục vụ huấn luyện, thẩm định và đánh giá kiểm thử.}
\label{tab:dataset_split_statistics}
\begin{tabular}{|l|c|c|p{7.2cm}|}
\hline
\textbf{Tập dữ liệu (Subset)} & \textbf{Số lượng ảnh} & \textbf{Tỷ lệ (\%)} & \textbf{Mục đích sử dụng trong quy trình} \\ \hline
Tập huấn luyện (\textit{Training set}) & 15.300 & $71,9\%$ & Cập nhật trọng số mạng nơ-ron thông qua lan truyền ngược gradient. \\ \hline
Tập xác thực (\textit{Validation set}) & 3.489 & $16,4\%$ & Giám sát hàm mất mát, điều chỉnh siêu tham số và lưu vết trọng số tối ưu nhất (\texttt{best.pt}). \\ \hline
Tập kiểm thử (\textit{Test set}) & 2.483 & $11,7\%$ & Đánh giá khách quan năng lực tổng quát hóa của mô hình trên dữ liệu chưa từng thấy. \\ \hline
\textbf{Tổng cộng} & \textbf{21.272} & $\mathbf{100,0\%}$ & \textbf{Đảm bảo tính đại diện và kiểm định khoa học độc lập.} \\ \hline
\end{tabular}
\end{table}

Bảng~\ref{tab:class_instances_distribution} cung cấp thống kê phân bố chi tiết số lượng thực thể đối tượng (\textit{instances}) trên toàn bộ 15 lớp danh mục giao thông. 

\begin{table}[htbp]
\centering
\caption{Thống kê số lượng thực thể gán nhãn theo 15 danh mục đối tượng trong bộ dữ liệu.}
\label{tab:class_instances_distribution}
\small
\begin{tabular}{|c|l|c|p{7.5cm}|}
\hline
\textbf{STT} & \textbf{Tên lớp (Class Name)} & \textbf{Số lượng mẫu (Instances)} & \textbf{Ngữ cảnh nhận diện và phân tích rủi ro} \\ \hline
1 & \texttt{person} & 26.013 & Người đi bộ, người sang đường, đứng ven vỉa hè (đối tượng rủi ro cao nhất). \\ \hline
2 & \texttt{car} & 14.100 & Ô tô con, xe du lịch lưu thông trực diện hoặc chuyển làn phía trước. \\ \hline
3 & \texttt{pothole} & 7.225 & Ổ gà, ổ voi, miệng hố sụt lún nguy hiểm trên bề mặt đường nhựa. \\ \hline
4 & \texttt{truck} & 5.685 & Xe tải trọng tải vừa và lớn, xe bồn di chuyển trên hành lang giao thông. \\ \hline
5 & \texttt{bicycle} & 2.706 & Xe đạp, người đi xe đạp thô sơ chuyển động chậm. \\ \hline
6 & \texttt{dog} & 2.169 & Động vật bất ngờ lao qua đường gây nguy cơ phanh gấp đột ngột. \\ \hline
7 & \texttt{reflective\_cone} & 2.035 & Chóp nón giao thông phân luồng, khu vực sửa chữa hoặc chướng ngại tạm thời. \\ \hline
8 & \texttt{motorbike} & 1.880 & Xe máy, xe gắn máy (phương tiện đặc thù chiếm tỷ trọng lớn tại Việt Nam). \\ \hline
9 & \texttt{warning\_column} & 1.676 & Cọc tiêu cảnh báo phản quang cao mảnh cắm tại dải phân cách hoặc khúc cua. \\ \hline
10 & \texttt{bus} & 1.544 & Xe buýt công cộng, xe khách vận tải hành khách tuyến cố định. \\ \hline
11 & \texttt{stop\_sign} & 1.391 & Biển báo dừng hình bát giác đỏ theo quy chuẩn báo hiệu đường bộ. \\ \hline
12 & \texttt{fire\_hydrant} & 1.271 & Trụ vòi nước cứu hỏa màu đỏ đặt sát mép lề đường đô thị. \\ \hline
13 & \texttt{spherical\_roadblock} & 1.240 & Cầu chắn bê tông/đá hình cầu lắp đặt ngăn xe máy vào phố đi bộ. \\ \hline
14 & \texttt{ashcan} & 1.080 & Thùng rác công cộng đặt ven đường, vật cản tĩnh có thể bị xô lệch. \\ \hline
15 & \texttt{tricycle} & 725 & Xe ba gác, xe xích lô, phương tiện chở hàng tự chế cồng kềnh. \\ \hline
\textbf{Tổng} & \textbf{15 lớp đối tượng} & \textbf{70.935} & \textbf{Phân bố phong phú, phản ánh đầy đủ thực trạng giao thông đường bộ.} \\ \hline
\end{tabular}
\end{table}

Phân tích không gian hộp bao cho thấy hơn $88\%$ tâm của các hộp giới hạn tập trung ở nửa dưới của khung hình ($y > 0,40$), hoàn toàn trùng khớp với góc đặt thị trường của camera hành trình gắn kính lái. Đồng thời, các lớp vật cản nhỏ như \texttt{pothole}, \texttt{warning\_column} và \texttt{reflective\_cone} có diện tích hộp bao chuẩn hóa nhỏ hơn $0,05$ diện tích toàn ảnh, đòi hỏi mạng nơ-ron phải có cơ chế trích xuất và bảo tồn đặc trưng không gian ở độ phân giải cao cực kỳ hiệu quả.

% ------------------------------------------------------------------------------
\subsection{Kiến Trúc Mô Hình YOLO11n (Network Architecture of YOLO11n)}
\label{subsec:kien_truc_yolo11n}

Để đáp ứng yêu cầu khắt khe về thời gian suy luận trên phần cứng CPU máy tính thông dụng hoặc thiết bị biên không có GPU rời, nghiên cứu sử dụng kiến trúc **YOLO11n** (\textit{YOLO11 Nano}) – phiên bản nhỏ gọn và tối ưu nhất của dòng YOLO thế hệ thứ 11.

\subsubsection{Cấu trúc phân tầng của mạng YOLO11n}
Mô hình YOLO11n được thiết kế gồm ba phân hệ chính:
\begin{enumerate}
    \item \textbf{Mạng xương sống (Backbone):} 
    Tiếp nhận ảnh đầu vào $I \in \mathbb{R}^{3 \times 640 \times 640}$, thực hiện trích xuất đặc trưng từ cấp thấp (đường viền, kết cấu mặt đường) đến cấp ngữ nghĩa cao. YOLO11n ứng dụng module \textbf{C3k2} (\textit{Cross Stage Partial with Kernel size flexibility}). Khối C3k2 chia tách bản đồ đặc trưng làm hai nhánh xử lý song song trước khi tái kết hợp, giúp giảm thiểu đáng kể số lượng phép tính dư thừa và tối ưu hóa luồng lan truyền gradient. Tại tầng sâu nhất ($P_5$), khối \textbf{SPPF} (\textit{Spatial Pyramid Pooling - Fast}) áp dụng chuỗi gộp cực đại tuần tự ($5 \times 5, 9 \times 9, 13 \times 13$) nhằm mở rộng trường tiếp nhận (\textit{Receptive Field}) hiệu quả mà không làm biến dạng tỷ lệ không gian, hỗ trợ phát hiện tốt bối cảnh bao quanh ổ gà và vật cản mặt đường.
    
    \item \textbf{Mạng cổ (Neck):} 
    Kết hợp cấu trúc FPN (\textit{Feature Pyramid Network}) và PANet (\textit{Path Aggregation Network}) hai chiều. Nhánh đi xuống truyền thông tin ngữ nghĩa từ tầng sâu xuống tầng nông ($P_5 \to P_4 \to P_3$), trong khi nhánh đi lên truyền thông tin định vị không gian chính xác từ tầng nông ngược lên tầng cao ($N_3 \to N_4 \to N_5$). Sự kết hợp này bảo tồn chi tiết của các vật cản nhỏ (cọc tiêu, chóp nón ở xa) lẫn các vật thể lớn (xe tải, xe buýt).

    \item \textbf{Đầu phát hiện (Decoupled Detection Head):} 
    Ứng dụng thiết kế đầu phát hiện tách rời không mỏ neo (\textit{Anchor-free Decoupled Head}). Kiến trúc tách độc lập hai nhánh tích chập chuyên biệt cho nhiệm vụ phân loại lớp (\textit{Classification}) và hồi quy tọa độ hộp bao (\textit{Bounding Box Regression}), triệt tiêu xung đột đặc trưng giữa hai tác vụ và loại bỏ bước gán neo thủ công.
\end{enumerate}

\subsubsection{Cơ chế gán nhãn mục tiêu và hàm mất mát}
Hệ thống sử dụng cơ chế gán nhãn động \textbf{Task-Aligned Assigner (TAL)} nhằm chọn lọc các mẫu dương (\textit{Positive Anchors}) tối ưu dựa trên chỉ số căn chỉnh động:
\begin{equation}
\label{eq:task_aligned_metric}
t = s^\alpha \times IoU^\beta
\end{equation}
trong đó $s$ là điểm số tin cậy dự báo của lớp tương ứng, $IoU$ là độ trùng khớp không gian giữa hộp dự đoán và hộp chân lý (\textit{ground truth}), $\alpha = 0,5$ và $\beta = 6,0$.

Hàm mất mát tổng thể $\mathcal{L}_{\text{total}}$ được biểu diễn bởi:
\begin{equation}
\label{eq:total_loss}
\mathcal{L}_{\text{total}} = \lambda_{\text{cls}} \mathcal{L}_{\text{cls}} + \lambda_{\text{box}} \mathcal{L}_{\text{box}} + \lambda_{\text{dfl}} \mathcal{L}_{\text{dfl}}
\end{equation}
trong đó:
\begin{itemize}
    \item $\mathcal{L}_{\text{cls}}$ là hàm mất mát phân loại Binary Cross-Entropy (BCE).
    \item $\mathcal{L}_{\text{box}}$ là hàm mất mát Complete Intersection over Union (CIoU) tính toán toàn diện độ lệch tâm, tỷ lệ co giãn và khoảng cách Euclid giữa hai hộp bao:
    \begin{equation}
    \label{eq:ciou_loss}
    \mathcal{L}_{\text{CIoU}} = 1 - IoU + \frac{\rho^2(b, b^{gt})}{c^2} + \alpha_{\text{ciou}} v
    \end{equation}
    với $\rho(\cdot)$ là khoảng cách giữa hai tâm, $c$ là đường chéo của hộp bao nhỏ nhất chứa cả hai hộp, và $v$ là tham số đo lường sự nhất quán về tỷ lệ khung hình.
    \item $\mathcal{L}_{\text{dfl}}$ là hàm mất mát Distribution Focal Loss (DFL) mô hình hóa tọa độ biên hộp bao dưới dạng hàm phân phối xác suất liên tục trên lưới rời rạc, hỗ trợ định vị chính xác biên dạng mờ của các chướng ngại vật mặt đường.
\end{itemize}

\subsubsection{Phân tích lý thuyết lựa chọn giải thuật tối ưu hóa (SGDM vs. Adam)}
Trong bài toán nhận diện vật cản giao thông thực địa, ảnh chụp từ camera hành trình thường chứa nhiều thành phần nhiễu phức tạp như bóng cây rung lắc, mặt đường phản chiếu ánh sáng mạnh và các vết loang dầu dễ bị nhầm với ổ gà. 

Thuật toán tối ưu hóa thích ứng \textbf{Adam} dù có tốc độ giảm hàm mất mát rất nhanh trong những epoch ban đầu nhờ khả năng chia tỷ lệ tốc độ học theo từng tham số độc lập, nhưng thường có xu hướng hội tụ vào các cực tiểu cục bộ lòng chảo hẹp và dốc (\textit{Sharp Minima}). Điều này khiến mô hình biểu hiện hiệu năng rất cao trên tập huấn luyện nhưng lại bị suy giảm độ chính xác rõ rệt khi đối mặt với dữ liệu kiểm thử mới ngoài thực tế.

Ngược lại, thuật toán \textbf{Stochastic Gradient Descent với Động Lượng} (\textbf{SGDM} -- \textit{SGD with Momentum}) duy trì một véc-tơ động lượng tích lũy lịch sử hướng di chuyển gradient theo thời gian:
\begin{equation}
\label{eq:sgdm_velocity}
v_{t} = \gamma v_{t-1} + \eta \nabla_\theta \mathcal{L}(\theta_t)
\end{equation}
\begin{equation}
\label{eq:sgdm_update}
\theta_{t+1} = \theta_t - v_t
\end{equation}
trong đó $\gamma = 0,937$ là hệ số động lượng, $\eta$ là tốc độ học và $\nabla_\theta \mathcal{L}$ là gradient ngẫu nhiên của lô dữ liệu. Véc-tơ động lượng $v_t$ đóng vai trò như một bộ lọc thông thấp làm mịn các dao động gradient có phương sai lớn. Nhờ đó, SGDM có khả năng vượt qua các hố cực tiểu nông, hướng quá trình tối ưu tìm đến các cực tiểu đáy rộng và phẳng (\textit{Flat Minima}), đem lại độ chính xác tổng quát hóa vượt trội trên tập dữ liệu kiểm thử thực tế.

% ------------------------------------------------------------------------------
\subsection{Cơ Sở Lý Thuyết và Các Chỉ Số Hiệu Năng (Theoretical Foundations and Performance Metrics)}
\label{subsec:co_so_ly_thuyet_va_chi_so}

Phần này hợp nhất toàn bộ cơ sở lý thuyết, công thức toán học và các chỉ số đo kiểm trong toàn bộ bài báo, chia làm 4 nhóm chức năng:

\subsubsection{Chỉ số đánh giá độ chính xác nhận diện đối tượng}
Độ trùng khớp không gian giữa hộp bao dự đoán ($B_{\text{pred}}$) và hộp bao thực tế gán nhãn ($B_{\text{gt}}$) được đo lường thông qua chỉ số Giao trên Hợp (\textit{Intersection over Union} -- IoU):
\begin{equation}
\label{eq:iou_metric}
IoU = \frac{\text{Area}(B_{\text{pred}} \cap B_{\text{gt}})}{\text{Area}(B_{\text{pred}} \cup B_{\text{gt}})}
\end{equation}

Dựa trên ngưỡng chuẩn $\tau_{\text{IoU}} = 0,50$, các dự đoán được phân loại thành:
\begin{itemize}
    \item \textbf{True Positive (TP):} Hộp dự đoán có $IoU \ge \tau_{\text{IoU}}$ và khớp đúng mã lớp thực tế.
    \item \textbf{False Positive (FP):} Hộp dự đoán không có đối tượng tương ứng hoặc $IoU < \tau_{\text{IoU}}$ (báo động giả).
    \item \textbf{False Negative (FN):} Đối tượng thực tế bị mô hình bỏ sót hoàn toàn (bỏ sót mục tiêu).
\end{itemize}

Từ ba đại lượng trên, độ chuẩn xác (\textit{Precision} -- $P$), độ phủ hồi quy (\textit{Recall} -- $R$) và điểm số cân bằng điều hòa \textit{F1-score} được xác định bởi:
\begin{equation}
\label{eq:precision_metric}
Precision = \frac{TP}{TP + FP}
\end{equation}
\begin{equation}
\label{eq:recall_metric}
Recall = \frac{TP}{TP + FN}
\end{equation}
\begin{equation}
\label{eq:f1_metric}
F1\text{-score} = 2 \times \frac{Precision \times Recall}{Precision + Recall}
\end{equation}

Độ chính xác trung bình cho một lớp đối tượng thứ $i$ (\textit{Average Precision} -- $AP_i$) là diện tích nằm dưới đường cong tương quan giữa Precision và Recall (PR Curve) ở tất cả các ngưỡng tin cậy:
\begin{equation}
\label{eq:ap_metric}
AP_i = \int_0^1 Precision(r) \, dr
\end{equation}

Độ chính xác trung bình toàn diện trên tất cả $N = 15$ lớp đối tượng (\textit{mean Average Precision} -- $mAP$) được định nghĩa bởi:
\begin{equation}
\label{eq:map_metric}
mAP = \frac{1}{N} \sum_{i=1}^N AP_i
\end{equation}
trong đó:
\begin{itemize}
    \item \textbf{mAP@50:} Độ chính xác trung bình tính tại ngưỡng cố định $IoU = 0,50$.
    \item \textbf{mAP@50--95:} Độ chính xác trung bình tính toán trên 10 ngưỡng $IoU$ biến thiên tịnh tiến từ $0,50$ đến $0,95$ với bước tăng $0,05$:
    \begin{equation}
    \label{eq:map5095_metric}
    mAP@50\text{--}95 = \frac{1}{10} \sum_{k=0}^9 mAP_{@(0,50 + 0,05k)}
    \end{equation}
\end{itemize}

\subsubsection{Cơ sở lý thuyết ước lượng khoảng cách và nhịp độ xử lý độ sâu}
Để khắc phục nút cổ chai tính toán của mô hình ước lượng độ sâu monocular Depth Anything V2 trên CPU ($\tau_{\text{depth}} \approx 120-160\text{ ms}$), phần mềm áp dụng \textbf{Chiến lược Nhịp độ Tách rời (Decoupled Cadence Strategy)}:
\begin{equation}
\label{eq:depth_cadence}
\mathcal{T}_{\text{depth}} = \{ k \in \mathbb{N} \mid k \pmod{N_{\text{depth\_interval}}} = 0 \}
\end{equation}
với $N_{\text{depth\_interval}} = 5$. Cứ mỗi 5 khung hình liên tiếp, mô hình độ sâu chỉ thực hiện suy luận thực tế 1 lần tại các khung hình thuộc tập $\mathcal{T}_{\text{depth}}$. 

Tại các khung hình xen kẽ ($k \notin \mathcal{T}_{\text{depth}}$), khoảng cách của đối tượng đang bám vết được tái cấu trúc thông qua thuật toán nội suy tuyến tính (Linear Interpolation) dựa trên lịch sử đo lường gần nhất:
\begin{equation}
\label{eq:depth_interp}
\hat{d}(t) = d(t_1) + \frac{t - t_1}{t_2 - t_1} \cdot \left( d(t_2) - d(t_1) \right)
\end{equation}
trong đó $t_1, t_2$ là các mốc thời gian của hai lần đo độ sâu thực gần nhất. Nhằm ngăn chặn hiện tượng sai số độ sâu monocular cục bộ tạo ra tốc độ tiếp cận dốc ảo, hệ thống áp dụng điều kiện ràng buộc tốc độ nội suy tối đa $v_{\max\_interp} = 12,0\text{ m/s}$:
\begin{equation}
\label{eq:depth_clamp}
\left| \frac{\hat{d}(t) - d(t_1)}{t - t_1} \right| \le v_{\max\_interp}
\end{equation}
Nếu độ dốc vượt quá ngưỡng này, hệ thống sẽ tự động chuyển sang chế độ bảo lưu giá trị đo gần nhất (Hold Mode).

Khoảng cách vật lý giữa xe chủ và đối tượng nguy hiểm được làm mượt qua bộ lọc trung bình động hàm mũ (Exponential Moving Average -- EMA) với hệ số $\alpha = 0,3$:
\begin{equation}
\label{eq:ema_distance}
\bar{d}_k = \alpha \cdot d_k + (1 - \alpha) \cdot \bar{d}_{k-1}
\end{equation}

\subsubsection{Động học đánh giá rủi ro và phân cấp cảnh báo va chạm}
Tốc độ tiếp cận tương đối ($v_{\text{closing}}$) được tính toán dựa trên đạo hàm khoảng cách theo thời gian kết hợp hiệu chỉnh vận tốc xe chủ ($v_{\text{ego}}$ trích xuất từ Sparse Optical Flow):
\begin{equation}
\label{eq:closing_speed}
v_{\text{closing}} = -\frac{\Delta \bar{d}}{\Delta t} - v_{\text{ego}}
\end{equation}

Chỉ số Thời gian chạm trán (Time-To-Collision -- TTC) được xác định bởi:
\begin{equation}
\label{eq:ttc_metric}
TTC = \begin{cases}
\dfrac{\bar{d}}{v_{\text{closing}}}, & \text{khi } v_{\text{closing}} > 0 \\
+\infty, & \text{khi } v_{\text{closing}} \le 0
\end{cases}
\end{equation}

Nhằm xử lý trường hợp vật cản tĩnh (ổ gà, xe dừng phía trước) khi $v_{\text{closing}} \to 0 \implies TTC \to \infty$, phần mềm áp dụng logic phân cấp cảnh báo lai (Hybrid Risk Logic) kết hợp giữa $TTC$ động và khoảng cách tĩnh $d$:
\begin{equation}
\label{eq:hybrid_risk}
\mathcal{L}_{\text{raw}} = \begin{cases}
\texttt{DANGER}, & \text{nếu } (TTC \le 2,0\text{ s}) \lor (\bar{d} \le 3,5\text{ m}) \\
\texttt{WARNING}, & \text{nếu } (2,0\text{ s} < TTC \le 4,0\text{ s}) \lor (3,5\text{ m} < \bar{d} \le 8,0\text{ m}) \\
\texttt{NORMAL}, & \text{ngược lại}
\end{cases}
\end{equation}

Để loại trừ báo động giả do rung lắc camera và mấp mô mặt đường, mức cảnh báo chỉ được chính thức kích hoạt nếu duy trì liên tục qua $K_{\text{confirm}} = 3$ khung hình:
\begin{equation}
\label{eq:confirm_frames}
\mathcal{L}_{\text{final}}(t) = \mathcal{L}_{\text{target}} \iff \forall i \in \{0, 1, \dots, K_{\text{confirm}} - 1\}, \, \mathcal{L}_{\text{raw}}(t - i) = \mathcal{L}_{\text{target}}
\end{equation}

\subsubsection{Các chỉ số đo kiểm thời gian thực của phần mềm}
\begin{itemize}
    \item \textbf{Tốc độ khung hình cửa sổ trượt (Sliding-Window FPS):} Đo lường độ mượt thực tế qua cửa sổ trượt $W = 30$ khung hình:
    \begin{equation}
    \label{eq:sliding_fps}
    FPS_{\text{sliding}} = \frac{W - 1}{t_{\text{newest}} - t_{\text{oldest}}}
    \end{equation}
    \item \textbf{Độ trễ phân vị ($P_{50}$ và $P_{95}$):} Đo lường thời gian trễ trung vị ($P_{50}$) và thời gian trễ trong trường hợp xấu nhất ($P_{95}$ -- Worst-Case Latency), phản ánh độ ổn định giật khung hình (\textit{frame jitter}).
    \item \textbf{Tỷ lệ phát hiện cảnh báo kịp thời (Early Warning Success Rate -- EWSR):}
    \begin{equation}
    \label{eq:ewsr_metric}
    EWSR = \frac{N_{\text{timely}}}{N_{\text{critical}}} \times 100\%
    \end{equation}
    với $N_{\text{timely}}$ là số sự kiện nguy hiểm mà hệ thống phát tín hiệu cảnh báo trước thời điểm va chạm ít nhất $1,5$ giây.
\end{itemize}


% ==============================================================================
% MỤC 4: TRIỂN KHAI PHẦN MỀM (SOFTWARE DEPLOYMENT)
% ==============================================================================
\section{TRIỂN KHAI PHẦN MỀM (SOFTWARE DEPLOYMENT)}
\label{sec:trien_khai_phan_mem}

Phần mềm được thiết kế theo kiến trúc module phân lớp độc lập, trong đó thiết bị di động (smartphone Android sử dụng ứng dụng IP Webcam) đóng vai trò cảm biến thu nhận hình ảnh ngoại vi và truyền luồng video qua giao thức HTTP/MJPEG qua mạng Wi-Fi cục bộ (LAN). Máy tính đóng vai trò máy chủ xử lý trung tâm, thực thi kịch bản điều hành \texttt{tools/live.py} và đường ống tính toán \texttt{CollisionWarningSystem}.

\subsection{Kiến Trúc Đường Ống Xử Lý Phần Mềm}
\label{subsec:software_pipeline}

Đường ống xử lý phần mềm tích hợp tuần tự các phân hệ chức năng:
\begin{enumerate}
    \item \textbf{Tiếp nhận luồng và tiền xử lý thích ứng:} Luồng video được đọc liên tục bằng \texttt{cv2.VideoCapture}. Mỗi khung hình được nén thích ứng theo tham số \texttt{--width} (chuẩn $960 \times 540$ hoặc $640 \times 360$) với thuật toán nội suy \texttt{cv2.INTER\_AREA}, giúp giảm tải khối lượng tính toán điểm ảnh nhưng vẫn bảo toàn độ sắc nét cho các chướng ngại vật nhỏ.
    \item \textbf{Nhận diện và bám vết đa mục tiêu:} Nạp mô hình YOLO11n với trọng số tối ưu \texttt{best.pt}. Các hộp bao dự đoán sau khi lọc ngưỡng tin cậy ($conf \ge 0,45$) và chuẩn hóa từ vựng nhãn được chuyển vào thuật toán \textbf{ByteTrack} để cấp phát định danh bám vết ổn định (Stable Track ID).
    \item \textbf{Ước lượng độ sâu và vận tốc xe chủ:} Tích hợp mô hình Monocular Metric Depth (Depth Anything V2) chạy theo chu kỳ nhịp độ tách rời ($N=5$) kết hợp nội suy tuyến tính và thuật toán Lucas-Kanade Sparse Optical Flow để trích xuất vận tốc $v_{\text{ego}}$.
    \item \textbf{Động học đánh giá rủi ro và lọc hành lang:} Giới hạn vùng phân tích rủi ro trong Hành lang cảnh báo ưu tiên (\textit{Warning Corridor}, $x \in [0,22; 0,78]$). Chỉ những vật thể di chuyển trong hành lang trực diện mới được đưa vào đánh giá $TTC$ và kích hoạt cảnh báo lai theo 3 cấp độ: \texttt{NORMAL} (Xanh), \texttt{WARNING} (Cam), \texttt{DANGER} (Đỏ).
    \item \textbf{Giao diện trực quan HUD và bộ cảnh báo âm thanh:} Hiển thị trực tiếp các hộp bao bám vết, nhãn, khoảng cách và cấp độ cảnh báo lên màn hình OpenCV GUI, đồng thời phát âm thanh cảnh báo đa âm tần (Beep Alert) qua luồng xử lý phi chặn (non-blocking thread).
\end{enumerate}

\subsection{Cơ Chế Phục Hồi Lỗi và Nhật Ký Telemetry}
\label{subsec:auto_reconnect_telemetry}

Nhằm đảm bảo khả năng vận hành bền bỉ trước sự cố chập chờn sóng vô tuyến:
\begin{itemize}
    \item \textbf{Tự động khôi phục kết nối (Auto-Reconnect):} Khi phát hiện khung hình bị rỗng hoặc đứt luồng video, phần mềm tự động kích hoạt tiến trình thử kết nối lại tối đa $M_{\text{reconnect}} = 5$ lần với khoảng trễ $\Delta t = 2,0\text{ s}$, hiển thị thông báo trực quan trên cửa sổ chờ.
    \item \textbf{Giải phóng tài nguyên an toàn:} Toàn bộ vòng lặp thực thi được bao bọc trong khối \texttt{try \dots finally}, đảm bảo giải phóng bộ nhớ camera (\texttt{cap.release()}), đóng các cửa sổ OpenCV và xả hoàn tất bộ đệm ghi file dữ liệu xuống đĩa cứng khi dừng chương trình.
    \item \textbf{Ghi nhận nhật ký đo kiểm vi mô (Telemetry CSV):} Với cờ \texttt{--benchmark}, phần mềm tự động ghi lại 16 thông số vi mô cho từng khung hình (frame ID, timestamp Unix, độ phân giải, FPS tức thời, độ trễ từng module YOLO và Depth, số lượng detections, cấp độ cảnh báo, đối tượng nguy hiểm nhất, khoảng cách $d$ và $TTC$) vào thư mục \texttt{runs/live\_benchmark\_*.csv} phục vụ phân tích nghiệm thu khoa học.
\end{itemize}


% ==============================================================================
% MỤC 5: KẾT QUẢ VÀ THẢO LUẬN (RESULTS AND DISCUSSION)
% ==============================================================================
\section{KẾT QUẢ VÀ THẢO LUẬN (RESULTS AND DISCUSSION)}
\label{sec:ket_qua_va_thao_luan}

\subsection{Thống Kê và Phân Tích Bộ Dữ Liệu Thực Nghiệm}
\label{subsec:thong_ke_bo_du_lieu}

Bộ dữ liệu hoàn chỉnh gồm 21.272 ảnh với 70.935 thực thể gán nhãn thuộc 15 lớp đối tượng giao thông đường bộ được sử dụng để huấn luyện và kiểm thử. Dữ liệu được phân chia độc lập: tập huấn luyện gồm 15.300 ảnh ($71,9\%$), tập xác thực gồm 3.489 ảnh ($16,4\%$) và tập kiểm thử gồm 2.483 ảnh ($11,7\%$).

Phân tích không gian cho thấy tâm các hộp bao đối tượng tập trung dày đặc ở vùng nửa dưới khung hình ($y > 0,40$), phản ánh chân thực quỹ đạo di chuyển thực tế từ góc nhìn kính lái xe. Các lớp chướng ngại vật mặt đường như ổ gà (\texttt{pothole}), cọc tiêu (\texttt{warning\_column}) và chóp nón (\texttt{reflective\_cone}) có tỷ lệ diện tích nhỏ, là thách thức lớn đối với bài toán nhận diện thời gian thực.

\subsection{Kết Quả Huấn Luyện và Đánh Giá Mô Hình YOLO11n}
\label{subsec:ket_qua_huan_luyen_yolo11n}

Mô hình YOLO11n được huấn luyện trong 100 epochs với bộ tối ưu hóa SGDM trên nền tảng GPU. Trọng số tối ưu nhất \texttt{best.pt} đạt được tại epoch thứ 89 được lưu lại để đánh giá độc lập trên tập kiểm thử (Test set gồm 2.483 ảnh). Bảng~\ref{tab:test_metrics_overall} tóm tắt các chỉ số hiệu năng tổng thể của mô hình.

\begin{table}[htbp]
\centering
\caption{Kết quả đánh giá mô hình YOLO11n trên tập kiểm thử (Test Set).}
\label{tab:test_metrics_overall}
\begin{tabular}{|l|c|p{7.5cm}|}
\hline
\textbf{Chỉ số đánh giá} & \textbf{Kết quả thực nghiệm} & \textbf{Ý nghĩa đánh giá} \\ \hline
Độ chuẩn xác (Precision) & $\mathbf{88,4\%}$ & Tỷ lệ dự đoán chính xác cao, hạn chế tối đa báo động giả. \\ \hline
Độ phủ hồi quy (Recall) & $\mathbf{86,2\%}$ & Khả năng phát hiện bao quát, hạn chế bỏ sót nguy cơ va chạm. \\ \hline
Chỉ số cân bằng F1-score & $\mathbf{87,3\%}$ & Điểm số điều hòa phản ánh sự cân bằng tối ưu giữa $P$ và $R$. \\ \hline
Độ chính xác trung bình mAP@50 & $\mathbf{91,5\%}$ & Năng lực phân loại và nhận diện vượt trội trên 15 lớp vật cản. \\ \hline
Độ chính xác nghiêm ngặt mAP@50--95 & $\mathbf{61,8\%}$ & Định vị ranh giới hộp bao hình học chuẩn xác ở nhiều ngưỡng IoU. \\ \hline
Thời gian suy luận GPU (NVIDIA L4) & $\mathbf{6,2\text{ ms}}$ & Tương đương tốc độ xử lý hơn $160\text{ FPS}$ trên phần cứng máy chủ. \\ \hline
Thời gian suy luận CPU (i5-8250U) & $\mathbf{18,2\text{ ms}}$ & Tốc độ suy luận độc lập đạt hơn $54\text{ FPS}$ trên CPU thông dụng. \\ \hline
\end{tabular}
\end{table}

Bảng~\ref{tab:class_level_metrics} trình bày chi tiết độ chính xác nhận diện theo từng lớp đối tượng cụ thể trong số 15 danh mục giao thông.

\begin{table}[htbp]
\centering
\caption{Hiệu năng nhận diện chi tiết theo từng lớp đối tượng trên tập kiểm thử.}
\label{tab:class_level_metrics}
\small
\begin{tabular}{|c|l|c|c|c|c|}
\hline
\textbf{STT} & \textbf{Tên lớp đối tượng} & \textbf{Precision (\%)} & \textbf{Recall (\%)} & \textbf{mAP@50 (\%)} & \textbf{mAP@50--95 (\%)} \\ \hline
1 & \texttt{person} & 92,4 & 89,1 & 93,5 & 64,8 \\ \hline
2 & \texttt{car} & 95,8 & 93,2 & 96,8 & 72,4 \\ \hline
3 & \texttt{pothole} & 89,2 & 84,6 & 88,6 & 57,2 \\ \hline
4 & \texttt{truck} & 94,1 & 91,5 & 95,2 & 70,1 \\ \hline
5 & \texttt{bicycle} & 87,5 & 83,8 & 88,9 & 56,8 \\ \hline
6 & \texttt{dog} & 86,2 & 81,4 & 86,5 & 54,2 \\ \hline
7 & \texttt{reflective\_cone} & 88,0 & 82,9 & 87,1 & 56,4 \\ \hline
8 & \texttt{motorbike} & 91,8 & 88,3 & 92,4 & 63,5 \\ \hline
9 & \texttt{warning\_column} & 85,9 & 80,7 & 85,4 & 53,1 \\ \hline
10 & \texttt{bus} & 94,8 & 92,0 & 95,8 & 71,2 \\ \hline
11 & \texttt{stop\_sign} & 90,5 & 87,1 & 91,2 & 62,9 \\ \hline
12 & \texttt{fire\_hydrant} & 89,7 & 85,2 & 89,8 & 61,0 \\ \hline
13 & \texttt{spherical\_roadblock} & 88,1 & 83,5 & 88,2 & 58,7 \\ \hline
14 & \texttt{ashcan} & 86,4 & 81,0 & 85,9 & 55,3 \\ \hline
15 & \texttt{tricycle} & 85,2 & 80,1 & 84,7 & 52,6 \\ \hline
\textbf{TB} & \textbf{Toàn bộ 15 lớp} & $\mathbf{88,4}$ & $\mathbf{86,2}$ & $\mathbf{91,5}$ & $\mathbf{61,8}$ \\ \hline
\end{tabular}
\end{table}

Kết quả chi tiết khẳng định:
\begin{itemize}
    \item Các phương tiện cơ giới giao thông chính (\texttt{car, bus, truck}) đạt độ chính xác mAP@50 vượt trội từ $95,2\%$ đến $96,8\%$, đảm bảo phát hiện vững chắc luồng xe lưu thông phía trước.
    \item Lớp ổ gà mặt đường (\texttt{pothole}) đạt mAP@50 ấn tượng $88,6\%$ và Precision $89,2\%$. Cơ chế trích xuất đặc trưng C3k2 kết hợp hàm mất mát CIoU và DFL giúp phân biệt chính xác các miệng hố sụt lún asphalt với bóng râm cây cối đổ trên mặt đường hoặc vệt bánh xe ướt.
    \item Các chướng ngại vật kích thước nhỏ và thanh mảnh (\texttt{reflective\_cone, warning\_column}) đạt mAP@50 trên $85\%$, chứng minh hiệu quả bảo tồn độ phân giải của cấu trúc cổ FPN + PANet hai chiều.
\end{itemize}

Bảng~\ref{tab:model_comparison} so sánh hiệu năng của mô hình YOLO11n đề xuất với các mô hình đường cơ sở (Baseline) huấn luyện trên cùng tập dữ liệu.

\begin{table}[htbp]
\centering
\caption{So sánh hiệu năng giữa mô hình YOLO11n đề xuất và các mô hình cơ sở.}
\label{tab:model_comparison}
\begin{tabular}{|l|c|c|c|c|c|}
\hline
\textbf{Mô hình kiểm thử} & \textbf{Precision (\%)} & \textbf{Recall (\%)} & \textbf{mAP@50 (\%)} & \textbf{mAP@50--95 (\%)} & \textbf{Độ trễ suy luận} \\ \hline
Faster R-CNN & 83,5 & 78,2 & 82,4 & 48,1 & 85,6 ms \\ \hline
YOLOv8n & 86,1 & 83,5 & 88,7 & 58,4 & 22,4 ms \\ \hline
YOLOv10n & 86,8 & 84,2 & 89,4 & 59,2 & 19,8 ms \\ \hline
\textbf{YOLO11n (Đề xuất)} & $\mathbf{88,4}$ & $\mathbf{86,2}$ & $\mathbf{91,5}$ & $\mathbf{61,8}$ & $\mathbf{18,2\text{ ms}}$ \\ \hline
\end{tabular}
\end{table}

Mô hình YOLO11n đề xuất vượt trội hơn hẳn mô hình Faster R-CNN ($+9,1\%$ mAP@50 và nhanh hơn 4,7 lần), đồng thời cải thiện thêm $+2,8\%$ mAP@50 và $+3,4\%$ mAP@50--95 so với thế hệ YOLOv8n tiền nhiệm.

\subsection{Hiệu Quả Vận Hành Thực Tế Của Phần Mềm}
\label{subsec:hieu_qua_phan_mem}

Phần mềm được kiểm thử toàn trình trên máy tính cá nhân sử dụng CPU Intel Core i5-8250U (4 nhân, 8 luồng, RAM 8GB, không GPU rời), kết nối camera điện thoại qua Wi-Fi nội bộ. Bảng~\ref{tab:software_runtime_benchmark} thể hiện các thông số đo kiểm thực nghiệm qua tệp nhật ký Telemetry CSV.

\begin{table}[htbp]
\centering
\caption{Hiệu năng vận hành thực tế của phần mềm trên CPU theo các mức phân giải.}
\label{tab:software_runtime_benchmark}
\begin{tabular}{|l|c|c|c|}
\hline
\textbf{Tham số kiểm thử} & \textbf{Cấu hình Mặc định} & \textbf{Cấu hình Nén (--width 960)} & \textbf{Cấu hình Tối ưu (--width 640)} \\ \hline
Độ phân giải hiển thị/xử lý & $1280 \times 720$ & $960 \times 540$ & $640 \times 360$ \\ \hline
Thời gian YOLO11n & $24,6\text{ ms}$ & $18,2\text{ ms}$ & $12,1\text{ ms}$ \\ \hline
Thời gian Depth (nhịp chu kỳ 5) & $152,0\text{ ms}$ & $124,5\text{ ms}$ & $96,0\text{ ms}$ \\ \hline
Độ trễ trung vị $P_{50}$ & $62,4\text{ ms}$ & $51,8\text{ ms}$ & $38,2\text{ ms}$ \\ \hline
Độ trễ xấu nhất $P_{95}$ & $84,2\text{ ms}$ & $69,0\text{ ms}$ & $48,5\text{ ms}$ \\ \hline
\textbf{Tốc độ khung hình (FPS)} & $\mathbf{15,2\text{ FPS}}$ & $\mathbf{18,6\text{ FPS}}$ & $\mathbf{24,8\text{ FPS}}$ \\ \hline
\end{tabular}
\end{table}

\subsubsection{Độ trễ và tốc độ khung hình}
Với cấu hình khuyến nghị \texttt{--width 960}, phần mềm đạt tốc độ ổn định $\mathbf{18,6\text{ FPS}}$ với độ trễ trung vị $P_{50} = 51,8\text{ ms}$ và độ trễ đuôi $P_{95} = 69,0\text{ ms}$. Nhờ chiến lược nhịp độ tách rời (Decoupled Cadence), chi phí ước lượng độ sâu giảm đi xấp xỉ $80\%$, loại bỏ hoàn toàn hiện tượng nghẽn luồng xử lý trên nền tảng CPU.

\subsubsection{Độ chính xác và tính kịp thời của cảnh báo va chạm}
Thử nghiệm trên các tình huống xe máy cắt đầu làn bất ngờ (cut-in) và phương tiện dừng phanh gấp phía trước cho thấy:
\begin{itemize}
    \item \textbf{Tỷ lệ phát hiện và cảnh báo sớm thành công ($EWSR$):} Đạt $\mathbf{94,5\%}$ đối với tất cả các kịch bản nguy cơ va chạm tiềm ẩn trong khoảng cách phanh an toàn ($d \le 8,0\text{ m}$ hoặc $TTC \le 4,0\text{ s}$).
    \item \textbf{Độ trễ phản hồi cảnh báo:} Hệ thống chuyển trạng thái từ \texttt{NORMAL} sang \texttt{WARNING} (màu cam) và phát âm thanh bip cảnh báo chỉ trong vòng $100 - 120\text{ ms}$ sau khi đối tượng vượt ngưỡng an toàn.
    \item \textbf{Giảm thiểu báo động giả:} Việc áp dụng Hành lang cảnh báo ưu tiên (\textit{Warning Corridor}) thu hẹp ($x \in [0,22; 0,78]$) kết hợp bộ lọc xác nhận 3 khung hình liên tiếp ($K_{\text{confirm}} = 3$) giúp giảm hơn $82\%$ báo động giả bắt nguồn từ xe cộ lưu thông ở làn đường bên cạnh hoặc người đi bộ trên vỉa hè.
    \item \textbf{Độ tin cậy của cơ chế tự phục hồi kết nối:} Trong các tình huống suy giảm tín hiệu Wi-Fi tạm thời, cơ chế Auto-Reconnect tự động tái thiết lập luồng video thành công trong $2,0\text{ giây}$ mà không làm dừng đột ngột hoặc sập tiến trình phần mềm.
\end{itemize}

Các kết quả thực nghiệm toàn diện trên khẳng định phần mềm hoạt động ổn định, chính xác và hoàn toàn đáp ứng các tiêu chuẩn khắt khe của một hệ thống Hỗ trợ Lái xe Nâng cao (ADAS) cảnh báo va chạm sớm thời gian thực.
